\section{System moments, initial conditions and the practical panel choice}
\label{s5:dynamic-panel}
The difference equation removes an individual effect; it does not create a
strong instrument. This distinction matters for persistent loan balances or
capital ratios. Consider first the scalar model
\[
 y_{it}=\alpha_i+\rho y_{i,t-1}+v_{it},\qquad |\rho|<1,
 \quad \E[v_{it}\mid\alpha_i,y_{i1},v_{i2},\ldots,v_{i,t-1}]=0.
\]
Independent entities and fixed $T$ define the sampling experiment. For
$s\le t-2$, both $\E[y_{is}v_{it}]$ and $\E[y_{is}v_{i,t-1}]$ vanish, giving
$\E[y_{is}\Delta v_{it}]=0$. These are the differenced moments already
derived above \citep{arellano1991}. Their validity does not require the
initial observation to have the stationary distribution.

Level equations permit a different restriction. Write
$y_{i1}=\alpha_i/(1-\rho)+e_{i1}$. Then
\[
 \Delta y_{i2}=(\rho-1)e_{i1}+v_{i2},\qquad
 \E[\Delta y_{i2}\alpha_i]=0
 \quad\text{if }\E[e_{i1}\alpha_i]=0.
\]
The recursion
$\Delta y_{it}=\rho^{t-2}\Delta y_{i2}+
\sum_{j=0}^{t-3}\rho^j\Delta v_{i,t-j}$ propagates this zero covariance.
Consequently $\E[\Delta y_{i,t-1}(\alpha_i+v_{it})]=0$.
Combining these level moments with differenced moments yields system GMM.
Mean stationarity is sufficient for this construction; equality of initial
and limiting variances is not required. Conversely a portfolio assembled at
loan origination or after a balance-sheet repair may systematically start
above or below its entity-specific long-run level. The additional moment can
then fail even when the differenced moments remain valid. This reconstructs
the initial-condition argument in \citet[Sections 3.1--3.3, pp.107--109]{blundellbond2023}.
The archived text is their 2023 account, not the original 1998 article.

Adding $\beta x_{it}$ requires a separate classification for $x$.
Contemporaneous endogeneity excludes $x_{it}$ as an instrument; lagged levels
must precede both innovations in $\Delta v_{it}$. Predetermination and strict
exogeneity allow different lag sets. A level moment using $\Delta x_{i,t-1}$
also needs $\E[\Delta x_{i,t-1}\alpha_i]=0$ and orthogonality to $v_{it}$.
No initial-condition theorem for $y$ automatically supplies this condition
for a bank's chosen action. With common macro shocks, cross-entity
independence itself fails and the fixed-$T$, large-$N$ argument must change.

\subsection{Weighting, covariance and what a correction can repair}
Let $b_i-G_i\theta$ collect the entity's moments after any prespecified lag
restriction and compression, where $\theta=(\rho,\beta)'$. For a positive
weight $W$, the quadratic criterion has minimizer
\[
 \hat\theta_W=(\bar G'W\bar G)^{-1}\bar G'W\bar b,
 \quad
 \sqrt N(\hat\theta_W-\theta_0)=
 (G'WG)^{-1}G'W\,N^{-1/2}\sum_i(b_i-G_i\theta_0)+o_p(1).
\]
The expansion needs $\bar G\to G$, a moment CLT, a nonsingular $G'WG$,
and consistent $W$. Thus the covariance estimator is
\[
 \widehat{\Var}(\hat\theta_W)=N^{-1}
 (\bar G'W\bar G)^{-1}\bar G'W\hat\Omega W\bar G
 (\bar G'W\bar G)^{-1},
\]
where $\hat\Omega$ is the covariance of the entire entity moment vector,
not of individual entity-date rows. Two-step GMM takes
$W=\hat\Omega(\hat\theta_1)^{-1}$ using an initial consistent estimator.
The first-order limit omits uncertainty in that estimated weight because
the population moment is zero. The finite-sample derivative is nevertheless
nonzero:
\[
 d\hat\theta_2=(\bar G'W\bar G)^{-1}\bar G'(dW)
 (\bar b-\bar G\hat\theta_2),\qquad
 dW=-W(d\hat\Omega)W.
\]
This displays the term whose variation a finite-sample correction must
address. The 2023 primary account explains the Windmeijer correction and its
motivation \citep[Section 2.3, p.105]{blundellbond2023}. The transparent local
reference implements the first-order sandwich, \emph{not} that correction;
the correction's original paper has not been substituted by a metadata record.
A production two-step adapter needs a source-qualified implementation of the
correction or a separately justified complete-refit alternative. Neither can
make invalid moments valid or rescue a rank-deficient Jacobian.

Compression also changes a statistical problem. An unrestricted triangular
lag matrix has $(T-2)(T-1)/2$ instruments in the scalar differenced model.
The reference collapses moments over dates and keeps lags two and three,
yielding four differenced moments for $(y,x)$, or six after adding level
moments. Scaling moments by observed root mean squares changes the one-step
weight; it leaves the population zero restrictions unchanged. The code
therefore reports this estimator as scaled identity-weight GMM, not the
canonical homoskedastic optimal one-step estimator. Small objective values
with many instruments can reflect overfitting. An overidentification test
is neither proof of exclusion nor evidence of adequate instrument strength.

\subsection{A variance correction does not fix the target of invalid moments}
\label{rel32:gmm-variance-target}
\providecommand{\USD}[1]{\ensuremath{\mathrm{USD}\,#1}}
The follow-up by \citet[Sections 2--4]{rel32hkl} supplies a primary
comparison of the two corrections. The inspected text is the May 2020
preprint, whose title page is dated 1 June; the later journal version
and the Windmeijer original are distinct texts. The follow-up derives
both the uncertainty from estimating the first-stage weight and the
additional variation carried by a nonzero overidentified moment vector.
Under correct specification, those residual-dependent contributions
vanish from the leading asymptotic variance. When no parameter makes
all population moments zero, they can remain first-order terms.

Continue with $g_i(\theta)=b_i-G_i\theta$ and independent entities.
For clarity, take the initial weight to be the identity and distinguish
the uncentered second moment from the covariance:
\[
 \theta_1=\arg\min_\theta\|\E g_i(\theta)\|^2,\qquad
 \Omega_1=\E[g_i(\theta_1)g_i(\theta_1)'],\qquad W=\Omega_1^{-1},
\]
\[
 \theta_2=\arg\min_\theta
       \{\E g_i(\theta)\}'W\{\E g_i(\theta)\},\qquad
 r_j=\E g_i(\theta_j),\quad H_1=G'G,\quad H_2=G'WG.
\]
The weight in the second criterion is evaluated at the first-stage
limit and held fixed during minimization. In particular, this is not
continuous updating. The first-order conditions give $G'r_1=0$ and
$G'Wr_2=0$; they do not imply that either residual vector vanishes.
Assume fixed moment and parameter dimensions, unique interior minima,
full column rank of $G$, nonsingular $\Omega_1$, and the finite moments
needed for the joint central limit theorem. The source's Assumption 1
states fourth-moment conditions for the first-stage moments and
Jacobians, and a second-moment condition at the second-stage target.
These are substantive regularity requirements, not a guarantee for
weak or proliferating instruments.

A direct differentiation of the two first-order conditions gives an
interpretable form of the source's equations (29)--(40). The first-stage
influence is
\[
 \psi_{1i}=H_1^{-1}
 \{G'[g_i(\theta_1)-r_1]+(G_i-G)'r_1\}.
\]
Let $\Omega_{1,j}$ denote the derivative of
$\E[g_i(\theta)g_i(\theta)']$ in coordinate $j$, evaluated at
$\theta_1$. Define the columns of $D$ and the second-stage influence
with the first-stage argument held fixed by
\[
 D_{\cdot j}=-H_2^{-1}G'W\Omega_{1,j}Wr_2,
\]
\[
 \begin{split}
 \kappa_i=H_2^{-1}\bigl\{&(G_i-G)'Wr_2
                +G'W[g_i(\theta_2)-r_2]\\
                &-G'W[g_i(\theta_1)g_i(\theta_1)'-\Omega_1]Wr_2
                    \bigr\}.
 \end{split}
\]
Then $\psi_{2i}=\kappa_i+D\psi_{1i}$ and the asymptotic covariance
of $\sqrt N(\hat\theta_2-\theta_2)$ is
\[
 \begin{split}
 V_2={}&\Var(\kappa_i)+D\Var(\psi_{1i})D'\\
       &+D\Cov(\psi_{1i},\kappa_i)
          +\Cov(\kappa_i,\psi_{1i})D'.
 \end{split}
\]
The covariance of the estimator itself is $V_2/N$. Each term has a
separate origin: moment variation, the estimated Jacobian, the estimated
second-moment matrix and the first-stage fit. Under correct specification,
$r_1=r_2=0$, so the residual-dependent pieces and $D$ disappear from
this population influence. Their fitted counterparts need not vanish
in a finite overidentified sample. Keeping them leads to the corrected
formulas in the source. If the first-stage weight is also estimated,
its influence supplies a further term; the identity-weight display here
must not be reused unchanged for an arbitrary panel adapter.

\subsection{An exact misspecified law distinguishes uncertainty from validity}
Consider the following project construction, chosen for transparent
arithmetic rather than as a fitted bank model. Let $U$ and $V$ be independent
fair signs, and set
\[
 G_i=\begin{pmatrix}1\\1+U/2\end{pmatrix},\qquad
 b_i=\begin{pmatrix}1+U\\3+2V\end{pmatrix}.
\]
The two population moments are $1-\theta$ and $3-\theta$; no common
zero exists. The first-stage target is $\theta_1=2$. Its uncentered
second moment and the resulting second-stage target are
\[
 \Omega_1=\begin{pmatrix}2&-2\\-2&6\end{pmatrix},\qquad
 \theta_2=\frac53,\qquad D=\frac{11}{36}.
\]
In the order $(U,V)=(-1,-1),(-1,1),(1,-1),(1,1)$, exact substitution gives
\[
 \psi_1=\left(-\frac54,\frac34,-\frac34,\frac54\right),\qquad
 \psi_2=\left(-\frac{151}{144},\frac{65}{144},
                     -\frac{97}{144},\frac{61}{48}\right).
\]
Both have mean zero. The full second-stage asymptotic variance is
$5827/6912$, approximately $.8430$. Retaining only variation in the
moment vector while fixing the Jacobian and weight gives $193/324$,
approximately $.5957$; the efficient correctly specified formula
$H_2^{-1}$ gives $2/3$. Neither is the full variance for this law.
Independent differentiation of the estimator functional with respect to
all four cell probabilities reproduces the displayed influences.
This checks the calculation, not finite-sample coverage.

Centering the weight also changes the target under misspecification.
Here $r_1=(-1,1)'$, so replacing $\Omega_1$ by
$\Omega_1-r_1r_1'$ gives the different second-stage target $3/2$.
This is why the source's centered-weight alternative changes both
its derivative and its pseudo-true estimand
\citep[Section 4, Remark 2]{rel32hkl}. It is not merely a different
standard-error label on an unchanged population object.

The banking consequence precedes any choice of confidence interval.
Suppose an analyst attaches the illustrative contribution rule
$100\theta-160$ dollars per account to these coefficients. The one-step,
uncentered two-step and centered two-step targets would produce
\USD{40}, \USD{6.67} and $-\USD{10}$, respectively. These are three
values of an assumed reporting map. None is an identified intervention
value, because the construction supplies inconsistent moment restrictions
and no causal response model. A more accurate standard error around one
pseudo-true target cannot justify treating it as the bank's policy effect.
The valid conclusion is to revisit the information and exclusion argument,
or explicitly defend an approximation target and its economic use.

\subsection{What this comparison permits an implementation to claim}
The follow-up's Theorem 1 concerns fixed-dimensional, full-rank linear
GMM with the stated sampling and moment conditions. Its correction is
not a general weak-instrument or many-instrument remedy. The paper also
distinguishes iterated GMM from continuous updating and explains why the
same finite-sample argument does not automatically extend to nonlinear
moment models \citep[Sections 1, 4--6]{rel32hkl}. The score construction
of \citet{kleibergen2021}, discussed below, answers a different question:
its robustness to identification and moment misspecification should not
be conflated with this variance correction.

A detail in the selected preprint's proof requires care. After equation
(94), it infers that a fifth derivative of an inverse-weight functional
vanishes because the underlying second-moment matrix is quadratic.
Matrix inversion does not preserve polynomial degree. In the finite
construction above, put
\[
 Q(\phi)=\{G'\Omega(\phi)^{-1}G\}^{-1}
              G'\Omega(\phi)^{-1}r_2
       =-\frac{2(\phi-24)(\phi-2)}{3(\phi^2+4\phi+36)}.
\]
Although $\Omega(\phi)$ is quadratic, $Q^{(5)}(2)=-725/31104\ne0$.
This disproves that proof step, not the asymptotic covariance conclusion.
The local Taylor argument instead uses smooth inversion on a neighborhood
where the weight and curvature stay nonsingular. A bounded second
derivative and a root-$N$ first-stage error give an unscaled quadratic
remainder of order $O_p(N^{-1})$, hence $o_p(1)$ after multiplication
by $\sqrt N$. The present derivation relies on this local argument;
it does not certify every claim or the uninspected journal version.

For the panel comparison that follows, the historical point estimates,
first-order sandwiches and all adverse coverage results remain exactly
as reported. This primary comparison supports a distinct implementation
candidate: preserve entity-level moment vectors, expose the weight
convention and target, include the required derivative and cross-covariance
terms, and then qualify the complete estimator on a declared design.
No corrected native panel result has been produced here. Replacing the
reported standard errors without rerunning that design would misstate
the evidence; treating misspecification robustness as instrument validity
would misstate the scientific question.

\subsection{A paired comparison and a banking decision rule}
The synthetic short-panel experiment uses $N=800$, eight transitions after
an initial observation, and
$y_{it}=\alpha_i+\rho y_{i,t-1}+.6x_{it}+v_{it}$, with
$x_{it}=.5x_{i,t-1}+.4v_{it}+\eta_{it}$. A 400-period burn-in generates
the regular initial law. The initial-condition challenger then adds a
multiple of $\alpha_i$ to the initial balance; the serial-error challenger
changes the innovation law. The four methods use identical datasets.
Coverage below is joint for both coefficients, using two-coordinate
Bonferroni 95\% intervals across 400 repetitions.
\begin{center}\small
\begin{tabular}{lrrrr}\toprule
Regime & Difference 1 & Difference 2 & System 1 & System 2\\\midrule
Stationary & .955 & .950 & .958 & .948\\
Persistent & .985 & .983 & .968 & .933\\
Initial departure & .950 & .960 & .033 & .000\\
Serial-error violation & .000 & .000 & .000 & .000\\\bottomrule
\end{tabular}\end{center}
One and two denote weighting steps, not different causal targets. Near
95\% coverage, the Monte Carlo standard error is about .011. The very high
coverage of a weak differenced design is not a precision advantage; weak
moments can produce broad intervals and unreliable normal approximations.
The initial-departure failure is the decisive exclusion from an automatic
system-GMM default. For established comparable entities, system moments can
be a challenger after a substantive initial-law argument. For new portfolios,
the primary differenced route has the weaker initial-law requirement, but
must still report Jacobian strength, instrument count and residual serial
correlation. If those fail, the result is an unsupported inferential claim,
not a prompt to search until a small standard error appears.

Misspecification-robust moment inference is another distinct alternative.
\citet[Sections 2--3, Assumptions 1--2]{kleibergen2021} targets the minimizer
of a continuously updated population GMM criterion when no parameter makes
all moments zero. Its joint moment/Jacobian CLT and covariance conditions
support an identification- and misspecification-robust score construction.
That pseudo-true minimizer need not equal a causal response or a meaningful
risk premium. This use of ``double robust'' concerns weak identification and
moment misspecification; it is different from outcome/propensity robustness
in Chapter~\ref{r3:learning}. The inspected copy is the 2021 preprint of work
identified in 2025 journal metadata. It is a substantive inference challenger,
not a license to keep a causal label after the identifying restrictions fail.
